\chapter{Conclusion}
\label{ch:conclusion}

This project produced a complete, safety-instrumented teleoperation stack for
a backpack-mounted pair of seven-degree-of-freedom arms commanded through a
miniature mannequin master, together with a characterisation of what that
configuration can and cannot do.

Against the objectives of \cref{ch:intro}: the master was designed and
fabricated and its sensing characterised, including a dated record of
progressive channel degradation; the mapping was derived and implemented,
including the coordinate-convention layer that isolates a class of error whose
failure mode is a near-full joint revolution; the safety architecture was
implemented and validated by fault injection rather than inspection; the
workspace was characterised, yielding the finding that determines the task
set; a verification methodology was established and applied across the
recorded runs; and the constraints blocking hardware validation were
identified by measurement rather than assumed.

The principal measured contributions are the workspace characterisation, with
its negative result and its quantified mechanical remedy; the account of what
degraded sensing costs and which two channels matter most; the finding that
the deployed clearance margin and the planner's collision test disagree across
the entire working band, resolvable laterally but not fore-and-aft; and the
observation that two independent constraints --- clearance and wrist
feasibility --- place the usable boundary in the same place.

The methodological contribution is the discipline of validating an instrument
against a known answer before reporting a surprising measurement taken with
it. This arose from necessity: across the project, unexpected results were
more often defects in measurement than in the system, and requiring every
analysis routine to recover a constructed answer caught a substantial number
of errors that would otherwise have entered the record.

\section{Future work}

The immediate steps are set by the limitations rather than by ambition, and
are ordered by how much each would improve the system.

\textbf{Repair the two identified master channels.} Regression shows the left
arm retains no reach observable without them, and every direct-teleoperation
baseline depends on them. This is the single highest-value action.

\textbf{Measure control-path latency.} The instrumentation is small --- seven
timestamps and a sequence number already partly present --- and it is the only
one of the planning report's targets that is both untouched and immediately
measurable. The result must be reported as a lower bound.

\textbf{Apply the mount change and re-derive.} The sweep quantifies the gain;
applying it invalidates the home pose, the workspace anchor and every
clearance figure, so it should be one deliberate pass with full
re-verification.

\textbf{Ingest and document the master hardware record.} The firmware, CAD and
bill of materials exist outside the repository; committing them would supply
the pin map, ADC resolution and fabrication parameters that
\cref{ch:method} currently cannot state.

\textbf{Then run the study.} The design is complete and costed at sixteen
dyads over a \SI{167}{\minute} session, comparing direct teleoperation against
shared autonomy across all five tasks.
\todo{MEASURE: the experimental programme in full --- no participant data
exists}
